\documentclass[11pt]{article}
\usepackage[margin=1in]{geometry}
\usepackage{setspace}
\doublespacing
\usepackage{parskip}
\usepackage{times}
\usepackage{titlesec}
\usepackage{fancyhdr}
\usepackage{xcolor}
\usepackage{amssymb}
\usepackage{enumitem}
\usepackage{titling}

\titleformat{\section}
  {\normalfont\Large\bfseries}{\thesection}{1em}{}
\titleformat{\subsection}
  {\normalfont\large\bfseries}{\thesubsection}{1em}{}
\titleformat{\subsubsection}
  {\normalfont\bfseries}{\thesubsubsection}{1em}{}

\pagestyle{fancy}
\fancyhf{}
\rhead{\textit{The Ninth Lock}}
\lhead{Page \thepage}
\cfoot{}

\setlength{\parindent}{0pt}

\begin{document}

\title{\Huge\bfseries The Ninth Lock}
\author{A Campaign for Fate's Edge\\\small\textit{The Overseer's Calculation}}
\date{}
\maketitle

\tableofcontents
\newpage

\section{Adventure Identity Statement}

\textbf{What makes this campaign sing?} A race against inevitability itself, where players must decide whether to help a desperate scholar uncover the truth of the world's oldest debt—or stop her before she opens a door that should remain sealed.

\textbf{What is the central tension?} Moira of Stonehaven, guided by a fragment of the demon Varash, believes she can end the cycle of the Ninth by opening the lost city of Khasurah. The players must decide: aid her, betray her, or save her from herself—while agents of the three sisters, rival factions, and the Hollow itself close in.

\textbf{What will players remember?} The weight of choice when the future is shown to be inevitable—yet not fixed. The cost of knowledge. The ninth cup un-poured.

\section{Thematic Pillars}

\begin{itemize}
    \item \textbf{Horror by Consequence:} The Ninth is not a monster; it is a debt that can never be fully paid. Horror comes from realizing that some doors, once opened, cannot be closed.
    \item \textbf{Mystery by Revelation:} The truth of Khasurah is revealed in fragments—each piece changes what the players think they know.
    \item \textbf{Exploration by Discovery:} The world-spanning quest for fragments reveals the layers of history beneath the Amaranthine.
\end{itemize}

\section{Campaign Structure}

\subsection{Act I: The Fragments (3–4 sessions)}

The players are recruited by Moira, a scholar turned ranger. She has located six fragments of the \textbf{Black Ledger of Khasurah}—scattered across the world. Her goal: assemble them to understand the Ninth taboo and perhaps find a way to end the Hollow's hunger.

\textbf{Opening Hook:} Moira meets the party in a neutral location (a Sidhi lighthouse, a Thepyrgos weigh-house, a Tulkani caravan camp). She shows them a fragment she already possesses—a torn page from the Black Ledger, the ink still wet though it is millennia old. ``Something is calculating. It has been calculating for nine thousand years. It will finish in nine more. We need to find the rest before it does.''

\subsubsection{The Fragments}

\begin{itemize}
    \item \textbf{The Ninth Line (Thepyrgos Archivolt):} Sealed in a lead box requiring eight witnesses to open. The ninth witness is always the Hollow. The party must find a way to open it without summoning it.
    \item \textbf{The Overseer's Ledger Page (Aeler Deep Vaults):} Behind a door with nine locks. The ninth key is a secret no living being remembers.
    \item \textbf{The Spider's First Cord (Crimson Valley):} Held by a Lethai-ar silk-weaver. She will only give it in exchange for witnessing a broken promise.
    \item \textbf{The Serpent's Shed Skin (Sekogo):} The Ukwe (Speaking Tree) demands an answer to a riddle: \textit{``What debt is never paid?''}
    \item \textbf{The Raven's Forgotten Name (Obishaal):} Hidden in Ikasha's dream. Requires entering the Ways Between and retrieving it from the sleeping goddess's mind.
    \item \textbf{Samir's Cut Page (Zakov):} In a sunken vault guarded by a drowned captain's ghost. He will only give it to someone who does not look back.
\end{itemize}

Each fragment retrieval is a mini-adventure (1 session). The order affects difficulty and information revealed.

\subsubsection{Timers (Visible to Players)}

\begin{itemize}
    \item \textbf{Moira's Certainty [8]:} Ticks when players help her, when she argues her case, or when the fragment's influence grows. At 4, she becomes more reckless. At 8, she is fully committed and cannot be reasoned with without Tam's intervention.
    \item \textbf{Sisters' Attention [6]:} Ticks when players uncover dangerous secrets, attract agents of Inaea, Isoka, or Ikasha. At 3, they are watched. At 6, an agent (Daughter of the Cord, Silk-Warden, dreamer) arrives to stop them.
    \item \textbf{Fragment's Whispers [4]:} Ticks when a character uses a fragment for knowledge or fails a Resolve test. At 4, a character gains the Condition \textit{Haunted by Varash} (the demon may offer a bargain).
    \item \textbf{Overseer's Calculation [10]:} Ticks automatically each session (or when players delay). When full, the Overseer completes the proof. The ninth lock opens if the party is not there to stop it.
\end{itemize}

\subsection{Act II: The Betrayal (2–3 sessions)}

After collecting three fragments, Moira reveals her true intent. She does not want to merely study Khasurah—she wants to \textit{open} it. The fragment of Varash has shown her a true future: the Overseer will complete its calculation in nine years, and the Clockwork Demon will awaken, scouring the world of inefficiency, debt, and mercy. She believes she can intervene.

\textbf{The players discover:} Moira has been lying—withholding information, manipulating their allegiances, even sabotaging rivals who might stop her. She is not evil; she is desperate. The fragment speaks in her dreams, showing her visions of a silent, perfect world without hunger.

\textbf{The Central Choice:}

\begin{enumerate}
    \item \textbf{Side with Moira:} Help her open Khasurah. Enter the drowned city (or its dream-echo) and face the Overseer.
    \item \textbf{Betray Moira:} Stop her from assembling the final fragments—steal them, alert the sisters' agents, or destroy the fragment (harming Moira in the process).
    \item \textbf{Save Moira:} Seek out Tam, her mentor, or another figure (the Gray Wanderer, a Daughter of the Cord) to intervene. Convince Moira to stop before it's too late.
\end{enumerate}

\subsubsection{Midpoint Crisis (Mandatory)}

Regardless of choice, the party is drawn into a threshold—either Khasurah (if siding with Moira) or a dream-echo of Khasurah (if trying to stop her). This is the final dungeon.

\subsection{Act III: The Threshold (2–3 sessions)}

The final location is \textbf{Khasurah's memory}—a psychic echo of the city's destruction, preserved by the Black Ledger fragments. It is unstable, shifting, and populated by Echoes of its citizens, Poltergeists of scholars, and Wardens left by the three sisters.

\textbf{Central Chamber:} A vast archive where the \textbf{Overseer's projection} sits at a desk, still calculating. The ninth lock is not a physical door; it is a \textit{mathematical proof}. If completed, it will open a gate for the Clockwork Demon.

\textbf{Moira's state:} If she is with the party, she is now fully under the fragment's influence. She attempts to complete the proof by inputting the final variable—the \textbf{name of Khasurah} (or the name of the original sin). Speaking the name unlocks the ninth lock.

\subsubsection{Final Choices (Multiple Outcomes)}

\begin{itemize}
    \item \textbf{Complete the proof with Moira:} The ninth lock opens. The Clockwork Demon begins to manifest. The campaign ends with a desperate escape (or a world changed irrevocably).
    \item \textbf{Destroy the fragments:} Shatter the Kal'rantha fragment or scatter the Black Ledger pages. Moira collapses, free but traumatized. The Overseer's calculation resets.
    \item \textbf{Replace the variable:} Change the proof—substitute a different name, a different condition. Seal the Clockwork Demon more permanently, or transfer the debt to a new entity.
    \item \textbf{Negotiate with the Overseer:} Convince the projection to stop—to accept that the variable cannot be known, the equation unsolvable. Requires a Spirit + Lore test (DV 6). The projection fades; the ninth lock remains closed forever.
\end{itemize}

\textbf{Tam's appearance:} If the players sought him in Act II, he arrives at the climax. He does not fight; he speaks to Moira: \textit{``You said you would return in a generation. It has not been a generation yet.''} Moira hesitates. The players act.

\section{Mechanical Integration}

\subsection{Core Innovation: The Fragment's Bargain}

\textbf{Signature System: Varash's Whisper}
\begin{itemize}
    \item \textbf{Purpose:} Represents the demon of Fate's influence—showing true futures to tempt characters into choices that serve inevitability.
    \item \textbf{Integration:} When a character touches a fragment or succeeds on a knowledge roll about Khasurah, the GM may offer a \textbf{Bargain}: accept a vision of a future outcome (true, but not inevitable) in exchange for a Story Beat or a tick on the \textit{Fragment's Whispers} timer.
    \item \textbf{Player Agency:} Players may refuse the bargain. Accepting gives them a genuine clue but accelerates their corruption.
    \item \textbf{Sample Uses:} \textit{``You see a vision of Moira standing before the Overseer, her hand on the final line. She is weeping. Do you accept this vision?''} (Tick Whispers +1, learn that Moira will not be able to stop herself.)
\end{itemize}

\subsection{Resource Management Layer}

\textbf{Khasuran Resonance:} A new tracker (0–3) representing how attuned the party is to the memories of the drowned city. Gain Resonance by:
\begin{itemize}
    \item Using a fragment to solve a problem (+1)
    \item Spending a Boon to recall a vision of Khasurah (+1)
    \item Accepting Varash's Whisper (+1)
\end{itemize}
At Resonance 3, the party can enter the dream-echo of Khasurah without an anchor—but they also attract the attention of the Wardens.

\textbf{Asset Building Opportunities:}
\begin{itemize}
    \item Each fragment, if not destroyed or returned, can become a \textbf{Minor Asset} (4 XP) granting +1 die to Lore or Investigation rolls concerning the Ninth.
    \item A character who successfully negotiates with the Overseer may gain a \textbf{Boon:} once per session, reroll a failed Resolve test when confronting inevitability (the `but you knew this would happen` effect).
\end{itemize}

\subsection{Timer Architecture}

See Act I for the four campaign timers. Summary:

\begin{itemize}
    \item \textbf{Moira's Certainty [8]} – Player actions and fragment influence.
    \item \textbf{Sisters' Attention [6]} – Player discoveries and dangerous actions.
    \item \textbf{Fragment's Whispers [4]} – Accepting visions, failing Resolve.
    \item \textbf{Overseer's Calculation [10]} – Automatic each session.
\end{itemize}

When a timer fills, immediate consequences occur. The GM should announce the tick and describe the effect.

\section{GM Support Systems}

\subsection{Session Preparation Checklist}
\begin{itemize}
    \item[$\square$] Advance \textit{Overseer's Calculation} by 1.
    \item[$\square$] Review which fragments remain and the party's relationship with Moira.
    \item[$\square$] Prepare 2-3 sensory details for each planned location.
    \item[$\square$] Note which timers are near filling; plan consequences.
    \item[$\square$] Have SB spend menus for the session's likely conflicts.
\end{itemize}

\subsection{Player Agency Reminders}

At choice points, ask:
\begin{itemize}
    \item ``What's the cost you're willing to pay?''
    \item ``Who benefits if you succeed here?''
    \item ``What would your character regret most?''
    \item ``How does this align with your drives?''
\end{itemize}

When tension lags:
\begin{itemize}
    \item Introduce a time pressure (a timer ticks audibly).
    \item Raise personal stakes (a character's contact is endangered).
    \item Complicate success (the fragment offers a truth they didn't seek).
    \item Force resource expenditure (a Wardens attack requires a Boon to avoid.)
\end{itemize}

\subsection{Complication Generator (SB Spend)}

\begin{itemize}
    \item \textbf{1 SB (Mild):} Environmental pressure – a floor tilts, a Warden shifts position.
    \item \textbf{2 SB (Moderate):} Tactical disadvantage – a rival faction arrives, a fragment's whisper reveals a false clue.
    \item \textbf{3 SB (Serious):} Major setback – a fragment is stolen, Moira's Certainty ticks +1.
    \item \textbf{4+ SB (Major):} Scene shift – the party is pulled into the dream-echo early, or a sister's agent demands a parley.
\end{itemize}

\section{Resolution Path Framework}

\subsection{Multiple Valid Approaches to the Climax}
\begin{enumerate}
    \item \textbf{Direct Confrontation:} Fight the Wardens and destroy the Overseer's projection.
    \item \textbf{Indirect Manipulation:} Use the fragments to rewrite the proof before Moira completes it.
    \item \textbf{Creative Solution:} Convince the Overseer that the variable is unknowable (negotiation).
    \item \textbf{Sacrificial Gambit:} One character offers their name as the final variable, sealing the lock with their own identity.
\end{enumerate}

\subsection{Outcome Matrix}
\begin{itemize}
    \item \textbf{Success:} The ninth lock remains closed. Moira is saved (or stopped). Cost: a character's memory, a fragment's destruction, or a new debt to the sisters.
    \item \textbf{Compromise:} The lock holds, but Moira is lost (she becomes a Hollowed, or wanders into Obishaal). The party gains a permanent enemy (a sister's faction).
    \item \textbf{Failure:} The ninth lock opens. The Clockwork Demon begins to emerge. The party must escape and find a new way to seal it (campaign continues).
    \item \textbf{Transformation:} The party chooses to replace the variable with something unexpected—a new covenant, a shared promise. The world changes in a small but profound way.
\end{itemize}

\section{Scalability and Modularity}

\subsection{Tier Adaptation}
\begin{itemize}
    \item \textbf{Lower Tiers (I–II):} Reduce \textit{Overseer's Calculation} to 6. Simplify fragment retrieval to 1-2 scenes each. Focus on Moira's personal arc; reduce sister involvement.
    \item \textbf{Higher Tiers (IV–V):} Add sub-plots for each character tied to their backgrounds. Expand sister agents into full factions. Allow players to travel to the original Khasurah (not just its echo).
\end{itemize}

\subsection{Session Modularity}
The campaign can be broken into:
\begin{itemize}
    \item \textbf{2-3 session arcs:} Acts I, II, and III as standalone mini-campaigns.
    \item \textbf{1-2 session adventures:} A single fragment retrieval (e.g., ``The Archivolt Heist'' or ``The Serpent's Skin'').
    \item \textbf{Standalone encounters:} A sister's agent confrontation, a vision from Varash, a side-quest to find Tam.
\end{itemize}

\section{Documentation Standards}

\subsection{Player-Facing Materials (Handouts)}
\begin{itemize}
    \item \textbf{Fragment Descriptions:} One-page summaries of each fragment's location, keeper, and rumor.
    \item \textbf{Timer Tracker:} A printable sheet with the four visible timers.
    \item \textbf{Moira's Journal Excerpts:} (from \textit{Up the Silverthread}) to establish her character and Tam's wisdom.
\end{itemize}

\subsection{GM Quick Reference}
\begin{itemize}
    \item \textbf{Fragment Keeper Motivations:} A table (Keeper, Desire, Weakness, Price).
    \item \textbf{Timer Triggers Summary:} One page.
    \item \textbf{SB Spend Menu for each Act:} Tailored to the threats.
    \item \textbf{The Overseer's Argument:} A script for the negotiation path.
\end{itemize}

\section{Quality Assurance Checklist}
\subsection*{Essential Elements}
\begin{itemize}
    \item[$\checkmark$] Clear adventure identity statement.
    \item[$\checkmark$] Meaningful player choices with consequences.
    \item[$\checkmark$] Integration with core Fate's Edge mechanics (Boons, SB, Timers, Position).
    \item[$\checkmark$] Custom mechanics (Varash's Whisper) that serve theme.
    \item[$\checkmark$] Multiple valid approaches to conflicts.
    \item[$\checkmark$] Clear escalation patterns (timers).
    \item[$\checkmark$] Memorable NPCs (Moira, Tam, fragment keepers).
    \item[$\checkmark$] Atmospheric location descriptions (see sample below).
    \item[$\checkmark$] Meaningful resource management (Resonance, fragments as Assets).
    \item[$\checkmark$] Consequences that ripple forward (faction attitudes, world changes).
\end{itemize}

\subsection*{Excellence Indicators}
\begin{itemize}
    \item[$\checkmark$] Mechanical innovation serves narrative (Varash's Whisper).
    \item[$\checkmark$] Player agency reinforced throughout (choice in Act II, multiple endings).
    \item[$\checkmark$] Thematic consistency maintained (debt, inevitability, choice).
    \item[$\checkmark$] Multiple resolution paths available.
    \item[$\checkmark$] Consequences feel earned and meaningful.
    \item[$\checkmark$] Scalability options clearly marked.
    \item[$\checkmark$] GM support materials comprehensive.
\end{itemize}

\section*{Sample Location: The Archivolt (Thepyrgos)}
\textit{The lead box sits on a pedestal of black basalt, its surface warm to the touch. Eight chains dangle from the ceiling, each ending in a brass handcuff. A ninth chain hangs empty, swaying gently in a breeze that has no source. The notary who led you here has already retreated to the door. ``Eight witnesses,'' she whispers. ``The ninth is always the Hollow. Do not let it be you.''}

\subsection*{The Ninth Line}
\textbf{Challenge:} Open the box without summoning the Hollow. \textbf{Solution:} Use a silver mirror to reflect the Hollow's attention. Requires an offering (a story, a memory) to distract it. \textbf{Partial:} The Hollow takes a minor memory (the name of a street, the face of a porter). \textbf{Failure:} The Hollow manifests and demands a name.

\section*{Epilogue: After the Campaign}
The world continues. The Overseer's calculation resets (or completes, depending on the ending). Moira either returns to her wanderings, settles in a village, or vanishes into the Ways Between. The players are changed. They have witnessed the oldest debt and chosen how to respond.

\vspace{2em}
\textit{The ink is still wet. The ninth page is blank. The Hollow reads over your shoulder.}

\end{document}
